\section{Introduction}
\label{sec:intro}

The modern machine learning serving paradigm is increasingly defined by multi-task specialization and on-the-fly model adaptation. With the advent of Parameter-Efficient Fine-Tuning (PEFT) techniques, such as Low-Rank Adaptation (LoRA) \cite{hu2021lora}, residual adapters \cite{rebuffi2017learning}, and prefix tuning \cite{li2021prefix}, practitioners can train lightweight, task-specific expert adapters. Under a dynamic, non-stationary test-time environment, serving a separate monolithic model or even routing sequentially to individual adapters is computationally and logistically prohibitive. Consequently, dynamic test-time model ensembling has emerged as a crucial research frontier \cite{ilharco2022editing, wortsman2022model, adamerging}, where task-specific expert weights are blended on-the-fly to construct a single, highly specialized, and adaptive network for each incoming query.

In real-world deployment scenarios, incoming queries do not arrive as independent and identically distributed (i.i.d.) batches. Instead, they manifest as continuous, sequential streams with significant temporal coherence (e.g., a user interacting with a model in a specific language, followed by a sudden pivot to a coding task). Navigating these non-stationary streams requires a \emph{stateful} routing mechanism that retains historical context to smooth out representation noise while remaining agile enough to pivot instantly under task transitions. 

\begin{figure}[t]
\centering
\includegraphics[width=0.9\linewidth]{comparison_plot.png}
\caption{Visualization of ensembling weight trajectories. Flat Euclidean approaches (such as Momentum-Merge and ChemMerge) rely on unconstrained linear updates followed by post-hoc Softmax normalization, warping the geometry and introducing significant lag and jitter. Our proposed Unitary Geodesic Routing (UGR) operates directly on the curved hypersphere $\mathbb{S}^{K-1}$, performing smooth geodesic rotations and mapping to the probability simplex via the square-root homeomorphism from Information Geometry, achieving near-zero jitter and zero representation lag under sudden transitions.}
\label{fig:teaser}
\end{figure}

To address this challenge, recent works have proposed stateful ensembling routers. For example, \emph{Momentum-Merge} implements a flat Exponential Moving Average (EMA) over expert activations, while \emph{ChemMerge} models router activations via continuous-time biochemical kinetics governed by ordinary differential equations (ODEs). While these methods outperform purely stateless routers like SABLE \cite{wang2020tent, liao2021are}, they suffer from a fundamental, shared methodological flaw: **they perform stateful updates in unconstrained flat Euclidean spaces.** 

Because the probability simplex $\Delta^{K-1}$ is a highly constrained, non-Euclidean manifold, flat updates in $\mathbb{R}^K$ must be projected post-hoc using operators like Softmax. This creates a severe disconnect between the optimization trajectory and the underlying probability space:
\begin{enumerate}
    \item \textbf{Representational Lag (Hysteresis):} Unconstrained linear updates accumulate inertia in flat space. Under a sudden task transition, the unconstrained state must "unwind" before the Softmax output reflects the switch, causing critical representational lag.
    \item \textbf{Geometric Distortion (Scale Mismatches):} Euclidean interpolation between probability vectors wiggles off the natural probability manifold, altering the norm and scale of intermediate feature activations and leading to mathematical instability.
    \item \textbf{High-Frequency Jitter:} To overcome lag, flat routers are forced to use high-temperature softmax or small inertia coefficients, which re-introduces high-frequency routing oscillations (jitter) in response to standard input noise.
\end{enumerate}

In this paper, we reject the flat-space assumption. We propose \textbf{Unitary Geodesic Routing (UGR)}, a non-Euclidean geometric approach that models ensembling states directly on the curved $(K-1)$-dimensional hypersphere $\mathbb{S}^{K-1} \subset \mathbb{R}^K$ (where ``Unitary'' refers strictly to the unit-norm constraint of our ensembling state vector). Rather than relying on Softmax projections, UGR leverages the square-root homeomorphism from classical Information Geometry (sharing a beautiful mathematical connection with Born's rule in quantum mechanics): $\alpha_{k,t} = (s_{k,t})^2$. This standard mapping homeomorphically relates the probability simplex $\Delta^{K-1}$ to the positive orthant of the unit hypersphere $\mathbb{S}^{K-1}_+$. Under this mapping, the celebrated Fisher-Rao Riemannian metric on the probability simplex corresponds exactly to the standard round Riemannian metric (geodesic arc-length) on the sphere. Thus, geodesic rotations on the hypersphere map to exact, closed-form Fisher-Rao geodesic flows on the probability simplex. This ensures that the ensembling weights and updates are mathematically guaranteed to lie on the probability simplex natively, yielding a Softmax-free state representation and projection pipeline with zero geometric distortion or scale-mismatch artifacts. While bottom-up target construction can still utilize localized Softmax gating for optimal task discrimination, our formulation natively supports a fully Softmax-free target construction alternative (e.g., using ReLU and $L_1$-normalization), providing a path toward entirely Softmax-free serving pipelines.

To model state transitions, we define a closed-form, Rodrigues-like geodesic rotation operator (Spherical EMA) that interpolates along the shortest great-circle path of the hypersphere. By doing so, we completely bypass expensive matrix exponentials and numerical ODE integrations, delivering extreme computational efficiency. To resolve the stability-plasticity dilemma, we introduce \emph{Torque-Driven Adaptive Agility}, which dynamically scales the angular step size in proportion to the angular distance (torque) between the current state and incoming activation signals. Finally, we formulate \emph{Spatial-Temporal Geodesic Coupling} to propagate the converged ensembling state smoothly across sequential boundary queries.

We evaluate UGR across both a high-fidelity synthetic 14-layer, 192-dimensional Analytical Coordinate Sandbox (ICS) and a multi-task text classification Mixture-of-Experts routing benchmark using the \texttt{20newsgroups} dataset (which acts as a highly controlled proxy for non-stationary text-routing, allowing us to mathematically isolate and analyze temporal router dynamics directly). Our experiments demonstrate that UGR establishes a new state-of-the-art Pareto frontier for test-time serving: in the synthetic sandbox, it achieves a Joint Mean Classification Accuracy of \textbf{75.08\%} ($\pm$ 1.14\%) while slashing layer-to-layer routing jitter ($L \ge 5$) by \textbf{2.10$\times$} compared to the continuous biochemical SOTA baseline ChemMerge; on the text classification stream, UGR delivers a Joint Mean Accuracy of \textbf{92.25\%} ($\pm$ 0.90\%), outperforming ChemMerge by a massive \textbf{+21.60\%} absolute margin while reducing routing jitter by \textbf{16.8$\times$}.

In summary, our key contributions are:
\begin{itemize}
    \item \textbf{Curved State-Space Formulation:} We establish the first test-time ensembling framework that operates entirely on the non-Euclidean hypersphere $\mathbb{S}^{K-1}$, mapping to the probability simplex natively via the information-geometric square-root homeomorphism (Fisher-Rao Geodesic Flow).
    \item \textbf{Closed-Form Geodesic Updates:} We derive a computationally efficient Rodrigues-like geodesic operator that performs spherical interpolation (Slerp) without numerical solvers.
    \item \textbf{Torque-Driven Agility:} We design a self-regulating control loop where angular torque dynamically scales the routing inertia, eliminating representational lag under sudden task switches.
    \item \textbf{Spatial-Temporal Coupling:} We couple sequential queries by propagating converged states across boundary conditions, ensuring temporal coherence.
    \item \textbf{Controlled Text-Routing Validation:} We validate UGR on a multi-task text classification Mixture-of-Experts pipeline, demonstrating its outstanding ability to filter out localized document noise and lock onto the correct task expert while acknowledging and discussing the scale gap toward full autoregressive LLM sequences.
\end{itemize}
